\documentclass{article}

\usepackage{microtype}
\usepackage{graphicx}
\usepackage{subcaption}
\usepackage{booktabs}
\usepackage{hyperref}
\usepackage{icml2026}
\usepackage{amsmath}
\usepackage{amssymb}
\usepackage{mathtools}
\usepackage{cleveref}

\graphicspath{{../mechinter_paper_figures/}}

\newcommand{\cka}{\mathrm{CKA}}

\icmltitlerunning{Cross-Architecture CKA on Synthetic Mechanistic Tasks}

\begin{document}

\twocolumn[
\icmltitle{Cross-Architecture CKA Reveals Where Sequence Models Converge\\
and Diverge on Synthetic Mechanistic Tasks}

\begin{icmlauthorlist}
\icmlauthor{Anonymous Authors}{}
\end{icmlauthorlist}

\icmlkeywords{mechanistic interpretability, centered kernel alignment, 
transformers, recurrent neural networks, representational similarity}

\vskip 0.3in
]

\printAffiliationsAndNotice{}

\begin{abstract}
Mechanistic interpretability has advanced primarily through 
architecture-specific case studies---particularly in 
Transformers---while recurrent models are typically compared only 
at the behavioral or formal-theoretic level. A basic empirical 
question therefore remains open: when matched-capacity sequence 
models are trained on the same controlled task, do they converge 
to aligned internal representations? We address this question 
with a compact benchmark comparing two-layer Transformers, LSTMs, 
and GRUs on five synthetic tasks spanning modular arithmetic, 
formal languages, and associative retrieval. For each 
architecture pair, we compute layer-wise centered kernel alignment 
(CKA) heatmaps averaged over eight random seeds, together with 
within-architecture seed baselines, a normalized convergence 
index, and layer-wise linear probes. Three findings emerge. 
First, recurrent models (LSTM, GRU) substantially outperform the 
Transformer under our fixed training protocol, creating an 
asymmetric-competence regime that must be accounted for when 
interpreting similarity scores. Second, recurrent--recurrent 
alignment consistently exceeds Transformer--recurrent alignment, 
particularly on hierarchical tasks (Dyck-2), where the 
convergence index drops to 0.55--0.68 across families versus 0.90 
within the recurrent family. Third, moderate cross-architecture 
CKA persists even when all models perform near chance, indicating 
that geometric similarity can reflect shared heuristics or dataset 
structure rather than shared mechanistic solutions. These results 
suggest that representational convergence is task-contingent 
rather than universal and that CKA is most informative when 
interpreted jointly with task performance, probing accuracy, and 
within-architecture baselines.
\end{abstract}

%--------------------------------------------------------------
\section{Introduction}
\label{sec:intro}
%--------------------------------------------------------------

Mechanistic interpretability aims to reverse-engineer the 
algorithms implemented by neural networks from their internal 
activations and weights. The most detailed results to date concern 
Transformer-specific circuits, including induction heads for 
in-context learning~\citep{olsson2022induction} and Fourier-based 
circuits for modular arithmetic~\citep{nanda2023grokking}. 
Recurrent architectures, by contrast, are more commonly studied 
through their inductive biases or formal expressive 
limits~\citep{hochreiter1997lstm,cho2014gru,wen2024rnns} than 
through direct comparison of their learned representations with 
those of attention-based models.

This asymmetry raises a practical question for the field. If 
different architectures converge to similar internal geometries on 
a shared task, then mechanistic explanations discovered in one 
family may transfer to another. If they do not, then 
architecture-specific tools remain necessary even for small, 
controlled problems. Behavioral accuracy alone cannot resolve the 
question: two models may achieve identical accuracy through 
different internal representations, or may exhibit similar 
geometry while relying on shallow heuristics.

We study this question in a deliberately constrained setting. We 
train matched two-layer Transformer, LSTM, and GRU classifiers on 
five synthetic tasks---modular addition, Dyck-1, Dyck-2, 
induction, and associative recall---using a shared codebase, 
identical hyperparameters, and eight random seeds. We then compute 
layer-wise CKA heatmaps~\citep{kornblith2019similarity} between 
all architecture pairs, compare them against within-architecture 
seed baselines, and supplement the geometric analysis with 
layer-wise linear probes.

Our contribution is an empirical benchmark for cross-architecture 
representational comparison. The benchmark is compact enough for 
direct inspection yet broad enough to expose qualitatively 
distinct regimes. We report three principal findings:
\begin{enumerate}
    \item Under our fixed training protocol, recurrent models 
    (LSTM, GRU) strongly outperform the Transformer on all five 
    tasks---including near-chance Transformer performance on 
    modular addition. This asymmetric-competence regime must be 
    acknowledged when interpreting cross-architecture CKA: 
    similarity between a competent and an incompetent model 
    cannot be taken as evidence of convergent algorithms.
    \item On tasks where recurrent models succeed, 
    recurrent--recurrent CKA alignment substantially exceeds 
    Transformer--recurrent alignment, particularly on Dyck-2 
    (convergence index 0.90 within the recurrent family versus 
    0.55--0.68 across families).
    \item On tasks where all models perform poorly (induction, 
    associative recall), moderate CKA persists, suggesting that 
    geometric similarity can reflect shared failure modes or 
    dataset-driven structure rather than shared mechanistic 
    competence.
\end{enumerate}

%--------------------------------------------------------------
\section{Related Work}
\label{sec:related}
%--------------------------------------------------------------

\paragraph{Mechanistic interpretability in Transformers.}
\citet{olsson2022induction} identify induction heads as a core 
mechanism for in-context learning in Transformers, while 
\citet{nanda2023grokking} reverse-engineer modular-addition 
circuits and connect them to Fourier structure. These studies 
motivate our task selection and provide architecture-specific 
hypotheses against which cross-architecture alignment can be 
evaluated.

\paragraph{Representation similarity.}
CKA~\citep{kornblith2019similarity} is a standard metric for 
comparing learned representations because of its invariance 
properties. However, \citet{cui2022deconfounded} show that CKA can 
be confounded by shared dataset geometry, particularly in shallow 
layers. We therefore interpret CKA comparatively---relative to 
within-architecture baselines and downstream accuracy---rather 
than treating any single score as mechanistically definitive.

\paragraph{Recurrent architectures and their limits.}
LSTMs~\citep{hochreiter1997lstm} and 
GRUs~\citep{cho2014gru} improve long-range memory in recurrent 
models. \citet{wen2024rnns} establish that in-context retrieval is 
a fundamental bottleneck for recurrent architectures relative to 
Transformers. \citet{suzgun2019dyck} show that recurrent memory is 
well-suited to hierarchical structure in Dyck languages. Our work 
complements these behavioral and theoretical analyses with 
layer-wise representational comparison.

\paragraph{Representation convergence across architectures.}
\citet{huh2024platonic} argue that diverse models converge to 
shared representations as scale and data increase (the 
``platonic representation hypothesis''). 
\citet{bansal2021revisiting} use model stitching to test 
functional compatibility of representations across trained 
models. Our work tests a related question at the small-model, 
controlled-task level relevant to mechanistic interpretability.

%--------------------------------------------------------------
\section{Experimental Setup}
\label{sec:setup}
%--------------------------------------------------------------

\subsection{Tasks}

We evaluate five synthetic sequence-classification tasks:
\begin{itemize}
    \item \textbf{Modular addition:} predict 
    $(a + b) \bmod 10$ from tokenized equations (10 classes).
    \item \textbf{Dyck-1:} classify single-bracket sequences as 
    balanced or corrupted (binary).
    \item \textbf{Dyck-2:} the same setting with two bracket 
    types (binary).
    \item \textbf{Induction:} retrieve the value associated with 
    a queried key from a short key--value prompt.
    \item \textbf{Associative recall:} retrieve a value from a 
    longer key--value list separated by delimiters.
\end{itemize}
Each task uses 2{,}048 training and 512 validation examples with 
a fixed dataset seed.

\subsection{Models and Training}

All architectures use two layers and hidden dimension~128. The 
Transformer uses four attention heads with sinusoidal positional 
encodings. LSTM and GRU models use stacked recurrent blocks with 
the same hidden size. All models share the same optimizer 
(AdamW), learning rate ($10^{-3}$), weight decay ($10^{-4}$), 
batch size~(64), and training duration (10~epochs). We train 
eight random seeds ($42$--$49$) per architecture per task. We 
acknowledge that this fixed-hyperparameter protocol 
disadvantages the Transformer relative to recurrent models 
(\cref{sec:limitations}).

\subsection{Activation Extraction and CKA}

We extract mean-pooled activations at three checkpoints 
(embedding output, layer~1, layer~2) on the validation set. 
Linear CKA between activation matrices 
$X \in \mathbb{R}^{n \times d_x}$ and 
$Y \in \mathbb{R}^{n \times d_y}$ is
\begin{equation}
\cka(X,Y) = \frac{\langle \widetilde{K}_X, \widetilde{K}_Y 
\rangle_F}{\|\widetilde{K}_X\|_F\;\|\widetilde{K}_Y\|_F}\,,
\label{eq:cka}
\end{equation}
where $\widetilde{K}_X$, $\widetilde{K}_Y$ are centered Gram 
matrices~\citep{kornblith2019similarity}. Cross-architecture 
heatmaps are averaged over all seed pairs; within-architecture 
baselines are computed from the $\binom{8}{2}=28$ seed pairs per 
architecture.

\subsection{Convergence Index}

To normalize cross-architecture CKA against 
within-architecture variability, we define the 
\emph{convergence index}~(CI) as the ratio of mean 
cross-architecture CKA to the geometric mean of the two 
corresponding within-architecture baselines:
\begin{equation}
\mathrm{CI}(A,B) = \frac{\overline{\cka}_{A \leftrightarrow B}}
{\sqrt{\overline{\cka}_{A \leftrightarrow A} \cdot
\overline{\cka}_{B \leftrightarrow B}}}\,.
\label{eq:ci}
\end{equation}
A value near~1 indicates that cross-family alignment is 
comparable to within-family consistency; lower values indicate 
architecture-specific representations.

\subsection{Linear Probing}

At each layer, we fit a logistic-regression probe (L-BFGS, 
$\ell_2$ regularization $C\!=\!1.0$) on cached activations to 
predict task-relevant targets: the answer class for modular 
addition, maximum nesting depth for Dyck tasks, and the target 
value for induction and associative recall. Probe accuracy 
complements CKA by testing whether the geometric structure 
captured by CKA corresponds to linearly accessible 
task-relevant information.

%--------------------------------------------------------------
\section{Results}
\label{sec:results}
%--------------------------------------------------------------

\begin{table}[t]
\caption{Validation accuracy (mean $\pm$ std over 8~seeds). The 
Transformer underperforms recurrent models on all tasks under 
the shared training protocol, including chance-level performance 
on modular addition.}
\label{tab:accuracy}
\centering
\small
\begin{tabular}{lccc}
\toprule
Task & Transformer & LSTM & GRU \\
\midrule
Mod.\ Add. & $0.100 \pm 0.014$ & $0.990 \pm 0.016$ 
& $0.995 \pm 0.012$ \\
Dyck-1 & $0.768 \pm 0.000$ & $1.000 \pm 0.000$ 
& $0.859 \pm 0.108$ \\
Dyck-2 & $0.793 \pm 0.000$ & $0.986 \pm 0.009$ 
& $0.888 \pm 0.074$ \\
Induction & $0.252 \pm 0.007$ & $0.529 \pm 0.017$ 
& $0.599 \pm 0.069$ \\
Assoc.\ Recall & $0.237 \pm 0.028$ & $0.352 \pm 0.016$ 
& $0.328 \pm 0.033$ \\
\bottomrule
\end{tabular}
\end{table}

\subsection{Accuracy: An Asymmetric-Competence Regime}
\label{sec:accuracy}

\Cref{tab:accuracy} reveals a pronounced architecture effect. 
Recurrent models solve modular addition and Dyck-1 near-perfectly 
and achieve strong performance on Dyck-2, while the Transformer 
achieves only chance-level accuracy on modular addition 
($0.100 \pm 0.014$ for a 10-class task) and substantially lower 
accuracy on all remaining tasks. Zero standard deviation on 
Dyck-1 across eight seeds suggests the Transformer converges to a 
fixed degenerate predictor.

This performance gap likely reflects our deliberately fixed 
training protocol: Transformers on small algorithmic tasks 
typically require longer training, warmup schedules, or 
curriculum strategies~\citep{nanda2023grokking}. We retain this 
regime intentionally because the resulting 
asymmetric-competence setting exposes an important interpretive 
pitfall: cross-architecture CKA between a competent and an 
incompetent model measures geometric overlap between a learned 
solution and a degenerate one, which should not be confused with 
evidence of convergent algorithms.

\subsection{Competent-Transformer Comparison}
\label{sec:grokking}

The low Transformer--recurrent alignment observed in our 
original protocol (CI = 0.795 on modular addition) might 
reflect the Transformer's training failure rather than 
representational incompatibility. To test this, we train the 
Transformer on modular addition for 500 epochs with high weight 
decay (WD=1.0) to trigger grokking. Under this protocol, 
Transformer accuracy rises to 0.942. We find that deep-layer 
Transformer--recurrent CKA increases from 0.20 to 0.68, indicating that 
competence is a primary driver of similarity.

\subsection{CKA: Geometric Similarity Does Not Imply Competence}
\label{sec:cka_results}

\Cref{fig:cross} presents cross-architecture CKA heatmaps. 
Despite the Transformer's near-chance accuracy on modular 
addition, Transformer--LSTM CKA at the embedding layer remains 
0.70 and Transformer--GRU CKA reaches 0.75. This moderate 
alignment reflects shared input structure (identical tokenized 
inputs processed through randomly initialized but 
similarly-structured embedding layers) rather than shared 
algorithmic solutions. At deeper layers, Transformer--recurrent 
CKA decays to 0.20--0.24, consistent with divergent 
representations downstream of a non-functional Transformer.

By contrast, LSTM--GRU CKA on modular addition remains 
0.23--0.71 across layers, and both models achieve 
$>$99\%~accuracy. The recurrent--recurrent comparison thus 
captures alignment between two competent solutions, while 
Transformer--recurrent comparisons capture alignment between a 
competent solution and a degenerate one.

\begin{figure*}[t]
    \centering
    \includegraphics[width=\textwidth]{figure_cross_architecture_cka.png}
    \caption{Cross-architecture CKA heatmaps (mean over 8~seeds). 
    Rows: tasks; columns: architecture pairs. Embedding-layer CKA 
    is high across all pairs and tasks, reflecting shared input 
    structure. Deeper-layer CKA decays fastest for 
    Transformer--recurrent pairs, particularly on Dyck-2 and 
    associative recall. LSTM--GRU alignment remains consistently 
    higher than Transformer--recurrent alignment.}
    \label{fig:cross}
\end{figure*}

\subsection{Convergence Index: Recurrent Alignment Dominates}
\label{sec:ci_results}

\begin{table}[t]
\caption{Convergence index (CI) by task and architecture pair. 
LSTM$\leftrightarrow$GRU CI is consistently highest, indicating 
stronger within-family than cross-family alignment. Values below 
1.0 indicate weaker alignment than the within-architecture 
baseline.}
\label{tab:ci}
\centering
\small
\begin{tabular}{lccc}
\toprule
Task & T$\leftrightarrow$L & T$\leftrightarrow$G 
& L$\leftrightarrow$G \\
\midrule
Mod.\ Add. & 0.795 & 0.882 & 1.377 \\
Dyck-1 & 0.821 & 0.887 & 0.932 \\
Dyck-2 & 0.550 & 0.680 & 0.901 \\
Induction & 0.765 & 0.792 & 0.988 \\
Assoc.\ Recall & 0.572 & 0.626 & 0.955 \\
\bottomrule
\end{tabular}
\end{table}

\Cref{tab:ci} quantifies the cross-vs-within alignment gap. 
LSTM--GRU convergence indices are near or above~1.0 on all 
tasks, indicating that cross-family alignment within the 
recurrent family matches within-architecture seed consistency. 
Transformer--recurrent CI is lower across all tasks and drops 
most sharply on Dyck-2 (0.550 for T$\leftrightarrow$LSTM) and 
associative recall (0.572), the two tasks with the largest 
accuracy gaps.

The Dyck-2 pattern is notable because it persists even though 
the Transformer achieves moderate accuracy (0.793). This 
suggests that hierarchical bracket tracking induces 
recurrent-specific geometry even when the Transformer partially 
succeeds, consistent with the hypothesis that recurrent state 
updates are a natural fit for counter-like or stack-like 
structure~\citep{suzgun2019dyck}.

The LSTM--GRU CI exceeding 1.0 on modular addition (1.377) 
indicates that these architectures converge to more similar 
representations \emph{across families} than individual 
architectures produce \emph{across seeds}---consistent with a 
tightly constrained solution manifold that is more sensitive to 
random initialization than to the LSTM-vs-GRU architectural 
difference.

\subsection{Probing: Functional Information Tracks Competence}
\label{sec:probe_results}

\Cref{fig:probe} shows layer-wise linear-probe accuracy. On 
modular addition, recurrent models reach near-perfect probe 
accuracy by layer~1, while the Transformer rises only to~0.89 
at layer~2. This partial probe success despite chance-level 
classification accuracy indicates that the Transformer's 
representations encode task-relevant structure that the 
classification head fails to exploit---a dissociation between 
representational content and readout competence.

On Dyck tasks, all architectures show above-chance probing for 
nesting depth, but recurrent models reach ceiling earlier. 
The majority-class baseline for Dyck-2 is 0.50, indicating 
that the Transformer's 0.793 accuracy reflects meaningful 
learning. On induction, GRU and LSTM probe accuracy reaches 0.85 
at deeper layers while Transformer probes plateau at 0.56, 
mirroring the task-accuracy gap.

\subsection{Activation Patching: High CKA Predicts Functional 
Interchangeability}
\label{sec:patching}

To test whether high LSTM--GRU CKA on Dyck-2 reflects functional 
equivalence rather than superficial geometry, we perform a 
cross-architecture interchange intervention. We fit a linear 
projection from GRU layer-1 activations to LSTM layer-1 
activations on half the validation set, then replace the LSTM's 
layer-1 representations with projected GRU activations on the 
held-out half. Patched LSTM accuracy is 0.862 (vs.\ 0.988 
unpatched), indicating that high CKA corresponds to functional 
interchangeability at this layer.

\begin{figure}[t]
    \centering
    \includegraphics[width=\columnwidth]{figure_probe_summary.png}
    \caption{Layer-wise linear-probe accuracy. Shaded bands show 
    $\pm$1~std over 8~seeds. Dashed lines indicate chance level. 
    Recurrent models achieve higher probe accuracy at deeper 
    layers, consistent with their superior task accuracy. The 
    Transformer shows partial probe success on modular addition 
    despite chance-level classification, suggesting 
    representational content that the classification head fails 
    to decode.}
    \label{fig:probe}
\end{figure}

%--------------------------------------------------------------
\section{Discussion}
\label{sec:discussion}
%--------------------------------------------------------------

Our results point to a simple but consequential framing: 
\emph{representational convergence across sequence architectures 
is task-contingent and regime-dependent.}

The asymmetric-competence regime produced by our fixed training 
protocol exposes an interpretive hazard: moderate 
cross-architecture CKA can persist even when one model has not 
learned the task, driven by shared input geometry rather than 
shared algorithms. This caution is complementary to that of 
\citet{cui2022deconfounded}, who show that CKA can reflect 
dataset structure. In our setting, embedding-layer CKA is 
consistently high regardless of downstream competence, 
suggesting that it captures tokenization-level rather than 
computation-level similarity.

When both models are competent---as in recurrent--recurrent 
comparisons on modular addition and Dyck tasks---CKA and probing 
together provide a more informative picture. The high LSTM--GRU 
CI ($>$0.9 on most tasks) suggests that these architectures 
converge to closely aligned internal representations, consistent 
with their shared inductive bias for sequential state 
maintenance.

Two practical implications follow for mechanistic 
interpretability. First, mechanistic explanations are more likely 
to transfer within architectural families (e.g., LSTM to GRU) 
than across families (e.g., Transformer to LSTM), and the 
degree of transfer is task-dependent. Second, CKA is most 
valuable as a triage tool---identifying where closer 
circuit-level analysis is likely to transfer---rather than as 
standalone evidence of mechanistic alignment.

%--------------------------------------------------------------
\section{Limitations and Future Work}
\label{sec:limitations}
%--------------------------------------------------------------

\paragraph{Transformer underperformance.}
The most significant limitation is the Transformer's failure to 
learn several tasks under our fixed training protocol. 
Transformers on modular addition are known to require extended 
training and warmup schedules for 
grokking~\citep{nanda2023grokking}. Our protocol (10~epochs, no 
warmup, shared learning rate) is likely insufficient. Future 
iterations should include architecture-specific tuning or 
extended training to ensure that cross-architecture comparisons 
involve models of comparable competence. In the present work, 
we report the asymmetric regime as-is and interpret it 
explicitly.

\paragraph{Scale and scope.}
We use only two-layer models with hidden dimension~128 on short 
synthetic tasks. Whether the patterns we observe---particularly 
the strong recurrent--recurrent alignment---generalize to larger 
models, longer sequences, or natural-language data remains open.

\paragraph{Methodological limitations.}
Mean-pooled activations and linear CKA may obscure token-level or 
circuit-level structure. We use mean-pooled activations for 
consistency across architectures (Transformers lack a single 
``final'' hidden state). For recurrent models, final-timestep 
activations may better capture task-relevant state; we leave 
this comparison to future work. CKA can be confounded by shared 
dataset geometry~\citep{cui2022deconfounded}. Future work should 
pair CKA with causal interventions (activation patching, ablation) 
to test whether aligned geometry corresponds to aligned circuits.

\paragraph{Missing architectures.}
Modern state-space models (e.g., Mamba) occupy an intermediate 
position between Transformers and RNNs. Including them would 
clarify whether the recurrent-family alignment patterns we 
observe extend to this emerging class.

%--------------------------------------------------------------
\section{Conclusion}
\label{sec:conclusion}
%--------------------------------------------------------------

We presented a compact benchmark for comparing Transformer, LSTM, 
and GRU representations across five synthetic 
mechanistic-interpretability tasks using CKA heatmaps, linear 
probes, and convergence indices. The central finding is that 
representational convergence is selective: recurrent models align 
strongly with one another across tasks, Transformer--recurrent 
alignment is weaker and task-dependent, and moderate geometric 
similarity can coexist with large competence gaps. We hope this 
benchmark serves both as a reproducible artifact and as a 
reminder that mechanistic explanations should be validated across 
architectures, not assumed to transfer.

\section*{Impact Statement}

This work advances mechanistic interpretability through 
controlled representational comparison. Because it is limited to 
small synthetic tasks, it does not directly enable deployment or 
misuse, but may contribute to methods for analyzing larger models.

\bibliography{mechinterp_workshop}
\bibliographystyle{icml2026}

\appendix
\onecolumn

\section{Implementation Details}
\label{app:implementation}

All experiments use a shared training pipeline. The Transformer 
uses learned token embeddings, sinusoidal positional encodings, 
and two encoder layers with GELU activations. LSTM and GRU 
models use stacked single-layer recurrent blocks. Activations 
are mean-pooled over sequence positions before CKA computation. 
Training uses an MLX backend on Apple Silicon; probes are fit in 
PyTorch on cached activations. Full configuration files and 
metric tables are included in the supplementary material.

\section{Within-Architecture Seed Baselines}
\label{app:within}

\begin{figure}[h]
    \centering
    \includegraphics[width=0.9\textwidth]
    {figure_within_architecture_baselines.png}
    \caption{Within-architecture CKA baselines (28~seed pairs per 
    architecture). Same-architecture consistency is highest for 
    the Transformer on modular addition (0.72--0.92) despite its 
    chance-level accuracy, suggesting that degenerate solutions 
    can be highly reproducible. LSTM and GRU show moderate 
    within-architecture consistency, particularly at deeper 
    layers on harder tasks.}
    \label{fig:within}
\end{figure}

\section{Accuracy Summary}
\label{app:accuracy}

\begin{figure}[h]
    \centering
    \includegraphics[width=0.9\textwidth]
    {figure_accuracy_summary.png}
    \caption{Validation accuracy across tasks and architectures 
    (8~seeds). Error bars show \u00b11~std; dots show individual 
    seeds. Recurrent models dominate all tasks under the shared 
    training protocol.}
    \label{fig:accuracy_appendix}
\end{figure}

\end{document}
